\documentclass{article}

% if you need to pass options to natbib, use, e.g.:
%     \PassOptionsToPackage{numbers, compress}{natbib}
% before loading neurips_2026

% The authors should use one of these tracks.
% Before accepting by the NeurIPS conference, select one of the options below.
% 0. "default" for submission
\usepackage{neurips_2026}
% the "default" option is equal to the "main" option, which is used for the Main Track with double-blind reviewing.
% 1. "main" option is used for the Main Track
%  \usepackage[main]{neurips_2026}
% 2. "position" option is used for the Position Paper Track
%  \usepackage[position]{neurips_2026}
% 3. "eandd" option is used for the Evaluations & Datasets Track
 % \usepackage[eandd]{neurips_2026}
 % if you need to opt-in for a single-blind submission in the E&D track:
 %\usepackage[eandd, nonanonymous]{neurips_2026}
% 4. "creativeai" option is used for the Creative AI Track
%  \usepackage[creativeai]{neurips_2026}
% 5. "sglblindworkshop" option is used for the Workshop with single-blind reviewing
 % \usepackage[sglblindworkshop]{neurips_2026}
% 6. "dblblindworkshop" option is used for the Workshop with double-blind reviewing
%  \usepackage[dblblindworkshop]{neurips_2026}

% After being accepted, the authors should add "final" behind the track to compile a camera-ready version.
% 1. Main Track
 % \usepackage[main, final]{neurips_2026}
% 2. Position Paper Track
%  \usepackage[position, final]{neurips_2026}
% 3. Evaluations & Datasets Track
 % \usepackage[eandd, final]{neurips_2026}
% 4. Creative AI Track
%  \usepackage[creativeai, final]{neurips_2026}
% 5. Workshop with single-blind reviewing
%  \usepackage[sglblindworkshop, final]{neurips_2026}
% 6. Workshop with double-blind reviewing
%  \usepackage[dblblindworkshop, final]{neurips_2026}
% Note. For the workshop paper template, both \title{} and \workshoptitle{} are required, with the former indicating the paper title shown in the title and the latter indicating the workshop title displayed in the footnote.
% For workshops (5., 6.), the authors should add the name of the workshop, "\workshoptitle" command is used to set the workshop title.
% \workshoptitle{WORKSHOP TITLE}

% "preprint" option is used for arXiv or other preprint submissions
 % \usepackage[preprint]{neurips_2026}

% to avoid loading the natbib package, add option nonatbib:
%    \usepackage[nonatbib]{neurips_2026}

\usepackage[utf8]{inputenc} % allow utf-8 input
\usepackage[T1]{fontenc}    % use 8-bit T1 fonts
\usepackage{hyperref}       % hyperlinks
\usepackage{url}            % simple URL typesetting
\usepackage{booktabs}       % professional-quality tables
\usepackage{amsmath}        % align, equation environments
\usepackage{amssymb}        % \mathbb, \mathcal
\usepackage{amsfonts}       % blackboard math symbols
\usepackage{nicefrac}       % compact symbols for 1/2, etc.
\usepackage{microtype}      % microtypography
\usepackage{xcolor}         % colors
\usepackage{algorithm}      % algorithm floats
\usepackage{algpseudocode}  % pseudocode
\usepackage{enumitem}       % compact lists

\newcommand{\soe}{\texttt{<SOE>}}
\newcommand{\eoe}{\texttt{<EOE>}}
\newcommand{\bos}{\texttt{<BOS>}}
\newcommand{\eos}{\texttt{<EOS>}}
\newcommand{\ttoken}[1]{\texttt{<#1>}}
\newcommand{\Hcal}{\mathcal{H}}

% Note. For the workshop paper template, both \title{} and \workshoptitle{} are required, with the former indicating the paper title shown in the title and the latter indicating the workshop title displayed in the footnote. 
\title{}


% The \author macro works with any number of authors. There are two commands
% used to separate the names and addresses of multiple authors: \And and \AND.
%
% Using \And between authors leaves it to LaTeX to determine where to break the
% lines. Using \AND forces a line break at that point. So, if LaTeX puts 3 of 4
% authors names on the first line, and the last on the second line, try using
% \AND instead of \And before the third author name.


\author{%
  % David S.~Hippocampus\thanks{Use footnote for providing further information
  %   about author (webpage, alternative address)---\emph{not} for acknowledging
  %   funding agencies.} \\
  % Department of Computer Science\\
  % Cranberry-Lemon University\\
  % Pittsburgh, PA 15213 \\
  % \texttt{hippo@cs.cranberry-lemon.edu} \\
  % examples of more authors
  % \And
  % Coauthor \\
  % Affiliation \\
  % Address \\
  % \texttt{email} \\
  % \AND
  % Coauthor \\
  % Affiliation \\
  % Address \\
  % \texttt{email} \\
  % \And
  % Coauthor \\
  % Affiliation \\
  % Address \\
  % \texttt{email} \\
  % \And
  % Coauthor \\
  % Affiliation \\
  % Address \\
  % \texttt{email} \\
}


\begin{document}


\maketitle


\begin{abstract}

\end{abstract}



\section{Introduction}

\section{Related Work}

\section{Methodology}
\label{sec:method}

We treat each patient's longitudinal record as a realization of a generative, marked temporal point process (MTPP) and learn it by autoregressive token prediction with a language-model backbone. Section~\ref{sec:problem} formalizes the setting and defines the prediction tasks. Section~\ref{sec:data} lists the two clinical datasets and the per-event feature schema the model consumes. Sections~\ref{sec:tokenization}--\ref{sec:inference} describe the three model variants studied in this paper: all three share the same backbone, adaptation scheme, and training objective, and differ only in the granularity at which the record is decomposed into events. A single optional extension adds a Hawkes-style intensity head trained against an auxiliary negative log-likelihood.

\subsection{Problem Formulation}
\label{sec:problem}

Let $\mathcal{S} = \{(t_k, \mathbf{x}_k)\}_{k=1}^{K}$ denote a patient record with $K$ observations, where $t_k \in \mathbb{R}_{\geq 0}$ is the time since baseline (months) and $\mathbf{x}_k$ is a heterogeneous, possibly partially observed mark comprising diagnostic labels, cognitive scores, neuroimaging summaries, and structured codes. Time-invariant covariates (sex, education, APOE-$\varepsilon$4 genotype, baseline age) are collected in a context vector $\mathbf{c}$. We write $\Hcal_k = \{(\mathbf{c}, (t_1, \mathbf{x}_1), \ldots, (t_k, \mathbf{x}_k))\}$ for the history up to and including the $k$-th observation and $\tau_k = t_k - t_{k-1}$ for the inter-observation interval.

A marked temporal point process is characterized by the conditional density of the next event given the history,
\begin{equation}
    p^*(t_{k+1}, \mathbf{x}_{k+1} \mid \Hcal_k)
    \;=\; \lambda^*(t_{k+1}, \mathbf{x}_{k+1}) \,
          \exp\!\Bigl(-\!\!\int_{t_k}^{t_{k+1}}\!\!\int \lambda^*(s, \mathbf{x})\,d\mathbf{x}\,ds\Bigr),
    \label{eq:mtpp}
\end{equation}
with conditional intensity $\lambda^*(\cdot \mid \Hcal_k)$. Rather than parameterizing $\lambda^*$ directly, we write each event as a finite string of discrete tokens and fit an autoregressive model $p_\theta(s_\ell \mid s_{<\ell})$ over the resulting sequence. The chain rule
$p^*(t_{k+1}, \mathbf{x}_{k+1} \mid \Hcal_k) = \prod_\ell p_\theta(s_\ell \mid s_{<\ell})$ then induces an implicit likelihood on the original $(t, \mathbf{x})$ events. This generative view unifies classification, regression, and forecasting under one loss.

\paragraph{Prediction tasks.} Given a history $\Hcal_k$ and a horizon $h>0$, we study five tasks: (T1) diagnosis at horizon $h$ (3-way classification); (T2) ADAS-Cog at $h$; (T3) a primary imaging biomarker at $h$ (ventricle-to-intracranial-volume ratio for TADPOLE; full-history class probability for EDUCATE); (T4) inter-visit interval $\tau_{k+1}$; and (T5) multi-step trajectory generation, in which $M$ future sequences are drawn by ancestral sampling. T1--T3 match established cross-sectional benchmarks; T4--T5 are enabled by the point-process formulation and have no cross-sectional counterpart.

\subsection{Data and Clinical Events}
\label{sec:data}

\paragraph{TADPOLE (ADNI).} The TADPOLE challenge subset of ADNI \cite{adni,marinescu2018tadpole} provides $\sim$1{,}700 patients in the training split (D1) with at least two visits each and a held-out rollover split (D2). Each visit yields up to eleven features: diagnosis $d \in \{\text{CN},\text{MCI},\text{AD}\}$, cognitive scores (MMSE, CDRSB, ADAS-13), and four MRI summaries (ventricles, hippocampus, whole brain, intracranial volume; ventricles and hippocampus are reported as ICV ratios). Four static attributes ($\mathbf{c}$) are sex, years of education, APOE-$\varepsilon$4 allele count, and baseline age.

\paragraph{EDUCATE.} EDUCATE is a tertiary-care electronic health record dataset covering dementia and related disorders, harmonized to OMOP concepts and MEDS events \cite{arnrich2024meds}. Each patient has an irregular stream of diagnosis, laboratory, medication, and procedure events with absolute timestamps. Targets are a binary healthy/non-healthy label (T1-binary) and a five-way class label (AD, MCI, FTD, DLB, other) restricted to non-healthy patients. Because diagnosis codes directly encode the label, we additionally evaluate with a \emph{leak-safe} preprocess that removes all dementia/ADRD ICD-10 prefixes (G30, G31, F00--F03, $\ldots$) and SNOMED/text mentions at load time; gaps between the two preprocess variants quantify code-level leakage.

Both datasets follow the same abstract schema $(t_k, \mathbf{x}_k)$ with $\mathbf{x}_k$ structured as modality-typed sub-vectors. This lets the three variants below apply uniformly; dataset-specific feature lists are deferred to Appendix.

\subsection{Three Granularity Variants}
\label{sec:tokenization}

All variants share a structural template for each patient:
\begin{equation}
    \bos \;\; \mathbf{c} \;\;
    \underbrace{\soe \; T(\tau_k) \; E(\mathbf{x}_k) \; \eoe}_{\text{event } k} \;\;
    \ldots \;\;
    \eos,
    \label{eq:template}
\end{equation}
and differ only in how the event payload $E(\mathbf{x}_k)$ is decomposed. The three variants trade off context length against the density of supervision (fraction of tokens under a cross-entropy loss whose prefix varies across visits). Let $F$ denote the number of observed feature dimensions at visit $k$ and $G$ the number of modality groups (imaging, cognitive, diagnosis, and, for EDUCATE, lab / medication / procedure).

\paragraph{(V) Visit-level encoding.} One event per visit. $E(\mathbf{x}_k)$ is a single row containing all observed feature--value pairs for that visit. This is the most compact representation: a 10-visit TADPOLE patient uses roughly 10 rows, fits comfortably in a 4{,}096-token window, and preserves all same-visit correlations at the cost of coarse per-event supervision.

\paragraph{(M) Modality-level encoding.} One event per modality group observed at the visit. $E(\mathbf{x}_k)$ is split into $G$ sub-rows sharing a timestamp, each carrying only the features of its group (e.g., an imaging sub-event with MRI summaries, a cognitive sub-event with MMSE/CDRSB/ADAS). Missing modalities are omitted rather than flagged, yielding a native handling of block-structured missingness.

\paragraph{(A) Atomic-event encoding.} One event per individual measurement. $E(\mathbf{x}_k)$ is decomposed into $F$ single-feature events of the form \ttoken{feat}$\;|\;t\;|\;v$. This gives the finest per-token supervision (one loss boundary per measurement) and the cleanest treatment of feature-wise missingness, at the cost of $F{-}1$ extra timestamp repetitions per visit.

Numerical values are emitted as text at the backbone tokenizer's native resolution; categorical labels (diagnosis, ICD/SNOMED codes) use a small set of added vocabulary items that we insert into the tokenizer and initialize from the mean of existing embeddings. Task prompts are appended after the history as $\ttoken{TASK\_*} \, \ttoken{horizon}$ so a single model answers all five tasks and multiple horizons from a shared backbone.

\subsection{Model and Adaptation}
\label{sec:model}

We use a causal language model $f_\theta$ as the backbone. Unless stated otherwise, $f_\theta$ is Qwen3.5 (0.8\,B and 2\,B variants studied), loaded in 4-bit NF4 quantization. For variants that introduce new tokens (category labels, modality tags, horizon bins), we extend the tokenizer and resize the embedding and output-projection matrices; new rows are initialized to the mean of the pretrained rows.

Trainable parameters are restricted to (i)~Low-Rank Adaptation \cite{hu2022lora} inserted in the attention and MLP projections ($\texttt{q,k,v,o\_proj}$, $\texttt{gate,up,down\_proj}$) as
$W = W_0 + (\alpha/r)\, B A$ with $r=16$, $\alpha=64$, dropout $0.1$, and (ii)~the new token rows and the language-model head. This reaches roughly $0.5\%$ of the backbone's parameters and fits a single 40\,GB GPU.

\paragraph{Optional intensity head.} As an orthogonal extension, we attach per-mark projections
$\mu_m, \alpha_m, \beta_m = \mathrm{softplus}\bigl(W_{\{\mu,\alpha,\beta\},m}\, h_k\bigr)$
from the backbone's last hidden state $h_k$ at event positions, defining a piecewise Hawkes intensity
\begin{equation}
    \lambda_m(t \mid \Hcal_k)
    \;=\; \mu_m + \!\!\sum_{j:\, t_j < t,\; m_j = m}\!\! (\alpha_m - \mu_m)\, e^{-\beta_m (t - t_j)}.
    \label{eq:hawkes}
\end{equation}
These heads are used to compute an auxiliary training loss (\S\ref{sec:training}) and to draw event times at inference (\S\ref{sec:inference}).

\subsection{Two-Stage Training}
\label{sec:training}

All three variants follow the same two-stage recipe, differing only in the token stream produced by their $E(\cdot)$.

\textbf{Stage~1 (continued pretraining).} For each patient with $K$ observations we emit every prefix of length $2,\ldots,K$ as an independent sequence (sub-sequence augmentation), which expands $\sim$1{,}700 TADPOLE patients into $\sim$10{,}000 sequences. The objective is next-token cross-entropy over the full sequence:
\begin{equation}
    \mathcal{L}_{\text{S1}}(\theta)
    \;=\; -\frac{1}{L} \sum_{\ell=1}^{L} \log p_\theta(s_\ell \mid s_{<\ell}).
    \label{eq:loss_s1}
\end{equation}

\textbf{Stage~2 (task supervised fine-tuning).} For each valid split point we construct prompt--response pairs $(P, R)$ where $P$ is the tokenized history followed by $\ttoken{TASK\_*}$ and a horizon token, and $R$ is the target. Examples across the five tasks are mixed. The loss is computed only on the response tokens:
\begin{equation}
    \mathcal{L}_{\text{S2}}(\theta)
    \;=\; -\frac{1}{|R|} \sum_{\ell=1}^{|R|} \log p_\theta(r_\ell \mid P, r_{<\ell}).
    \label{eq:loss_s2}
\end{equation}
We study two regimes: \emph{next-visit} ($R$ targets only the immediate next observation) and \emph{horizon-conditioned} ($R$ targets any future observation whose offset matches the horizon token in $P$, so one model answers arbitrary horizons).

\textbf{Optional auxiliary loss.} When the intensity head is attached, Stage~2 minimizes
$\mathcal{L} = \mathcal{L}_{\text{S2}} + \gamma\,\mathcal{L}_{\text{NLL}}$ with the standard MTPP negative log-likelihood
\begin{equation}
    \mathcal{L}_{\text{NLL}}
    \;=\; \int_0^{T} \sum_m \lambda_m(s)\, ds \;-\; \sum_{k=1}^{K} \log \lambda_{m_k}(t_k),
    \label{eq:nll}
\end{equation}
computed with a trapezoidal quadrature on the per-event intensities of Eq.~\eqref{eq:hawkes}. We linearly warm $\gamma$ from $0$ to $0.1$ over one epoch so that the text objective is established before the intensity head contributes gradients. Hyperparameters (learning rate $10^{-4}$ in both stages, effective batch size $16$, 5--9 epochs) are listed in Appendix.

\subsection{Inference}
\label{sec:inference}

\paragraph{Deterministic tasks (T1--T4).} We decode greedily at temperature $T_{\text{gen}}=0.05$ with constrained generation for categorical outputs. Scalar targets (ADAS, ventricles, inter-visit time) are read either as the generated numerical string or, for variants that emit a quantile token, as the expectation
$\hat{y} = \sum_j \mathrm{softmax}(z_{j})\,\mu_j$, where $z_j$ is the logit of the $j$-th bin token and $\mu_j$ the bin mean computed on the training split; the expectation is strictly more informative than the argmax bin.

\paragraph{Trajectory generation (T5).} We draw $M{=}100$ trajectories per patient by ancestral sampling at $T_{\text{gen}}=0.7$ conditioned on $\Hcal_k$ followed by $\ttoken{TASK\_TRAJ}$ (Alg.~\ref{alg:traj}). Each trajectory is parsed back into events; aggregates across $M$ samples yield event probabilities, expected biomarker curves with Monte-Carlo confidence bands, and a validity rate measuring how many trajectories have well-formed structure and clinically admissible transitions (no AD$\!\to\!$CN reversions).

\begin{algorithm}[t]
\small
\caption{Monte-Carlo trajectory generation}
\label{alg:traj}
\begin{algorithmic}[1]
\Require history $\Hcal_k$, sample count $M$, temperature $T_{\text{gen}}$, maximum length $L_{\max}$
\For{$m = 1, \ldots, M$}
    \State $\mathrm{seq} \leftarrow \Hcal_k \,\|\, \ttoken{TASK\_TRAJ}$
    \While{$|\mathrm{seq}| < L_{\max}$ \textbf{and} last token $\neq \eos$}
        \State $x \sim \mathrm{Cat}\bigl(\mathrm{softmax}(f_\theta(\mathrm{seq}) / T_{\text{gen}})\bigr)$
        \State $\mathrm{seq} \leftarrow \mathrm{seq} \,\|\, x$
    \EndWhile
    \State $\hat{\mathcal{S}}_m \leftarrow \mathrm{Parse}(\mathrm{seq})$
\EndFor
\end{algorithmic}
\end{algorithm}

\paragraph{Intensity-based time sampling.} When the intensity head is attached, inter-visit times may alternatively be drawn by inverse-CDF sampling on Eq.~\eqref{eq:hawkes}: evaluate $\Lambda(t) = \int_{t_k}^{t} \sum_m \lambda_m(s)\,ds$ on a grid, then draw $u \sim \mathcal{U}(0,1)$ and return $t^\star$ with $1 - e^{-\Lambda(t^\star)} = u$; the next mark is drawn from $\lambda_m(t^\star)/\sum_{m'} \lambda_{m'}(t^\star)$. This produces continuous-time samples and is compared against token-level sampling in experiments.

\section{Experiments}

\section{Discussion \& Conclusion}




\section*{References}



%%%%%%%%%%%%%%%%%%%%%%%%%%%%%%%%%%%%%%%%%%%%%%%%%%%%%%%%%%%%

\appendix



%%%%%%%%%%%%%%%%%%%%%%%%%%%%%%%%%%%%%%%%%%%%%%%%%%%%%%%%%%%%

\newpage
\input{checklist.tex}


\end{document}